\section{Leafwise cohomology}


    \subsection{Leafwise de Rham complex}
    
    For the triple $(M,\mathcal{F},\phi)$, let $T\mathcal{F}$ be a sub-bundle of the tangent bundle $TM$ which is tangent to the leaves of the foliation.
    The restriction of $T\mathcal{F}$ on a leaf $\mathcal{L}$ is identified with the tangent bundle $T\mathcal{L}$ of the leaf.
    We call $T\mathcal{F}$ the \textit{leafwise tangent bundle}.

    We define the space of leafwise $i$-forms by
    \begin{equation*}
        \Omega^{i}_{\mathcal{F}}(M):=\Gamma(M,\wedge^{i}T^{*}\mathcal{F})\subset{\Omega^{i}(M)}.
    \end{equation*}
    Let $d_{\mathcal{F}}$ (resp. $d_{0}$) be the exterior derivative acting only along leaves (resp. the flow).
    Then the de Rham complex $(\Omega^{i}(M),d^{i})$  has a decomposition:
    \begin{center}
    \begin{tikzcd}
    \cdots \ar[r]\ar[dr, ""] & \Omega^{i}_{\mathcal{F}}(M) \ar[d, "\oplus"]\ar[r, "d_{\mathcal{F}}^{i}"]\ar[dr, "d_{0}^{i}"] & \Omega^{i+1}_{\mathcal{F}}(M)  \ar[d, "\oplus"]\ar[r]\ar[dr, ""] & \cdots\\
    \cdots \ar[r] & \Omega^{i}_{0}(M) \ar[r, "d_{\mathcal{F}}^{i}"] & \Omega^{i+1}_{0}(M) \ar[r] & \cdots
    \end{tikzcd}
    \end{center}
    where $\Omega_{0}^{i}(M)$ is the complement of $\Omega^{i}_{\mathcal{F}}(M)$.
    
    We simply denote the restriction $d^{i}|_{\Omega^{i}_{\mathcal{F}}(M)}$ by $d_\mathcal{F}^{i}$.
    Since we have $d^{i+1}_{\mathcal{F}}\circ{d}^{i}_{\mathcal{F}}=0$ on $\Omega^{i}_{\mathcal{F}}(M)$, 
    the pairs $\{(\Omega^{i}_{\mathcal{F}}(M), d^{i}_{\mathcal{F}})\}_{i}$ form a cochain complex:
    \begin{equation*}
        0\overset{}{\rightarrow}\Omega^{0}_{\mathcal{F}}(M)\overset{d^{0}_{\mathcal{F}}}{\rightarrow}\Omega^{1}_{\mathcal{F}}(M)\overset{d^{1}_{\mathcal{F}}}{\rightarrow}\Omega^{2}_{\mathcal{F}}(M)\overset{d^{2}_{\mathcal{F}}}{\rightarrow}0.
    \end{equation*}
    We call the complex \textit{leafwise de Rham complex}.
    
    \subsection{Leafwise cohomology}
    We denote the kernel of $d_{\mathcal{F}}^{i}$ by $Z_{\mathcal{F}}^{i}(M)$ and the image of $d_{\mathcal{F}}^{i}$ by $B_{\mathcal{F}}^{i+1}(M)$.
    Note that a leafwise $i$-th form in $Z_{\mathcal{F}}^{i}(M)$ (resp. $B_{\mathcal{F}}^{i+1}(M)$) is called leafwise closed $i$-th form (resp. leafwise exact $i$-th form).
    
    \begin{definition}[Leafwise cohomology]
    We define the $i$-th leafwise cohomology group by
    \begin{equation*}
        H^{i}_{\mathcal{F}}(M):={Z_{\mathcal{F}}^{i}(M)}/{B_{\mathcal{F}}^{i}(M)}.
    \end{equation*}
    \end{definition}
    The leafwise cohomology group is trivial for $i>2$.
    
    Unfortunately, the leafwise cohomology group is of infinite dimension in general and not a Hausdorff space.
    We modify it by taking a quotient with respect to the closure in the smooth topology.
    \begin{equation*}
        \bar{H}^{i}_{\mathcal{F}}(M):= {Z_{\mathcal{F}}^{i}(M)}/{\overline{B_{\mathcal{F}}^{i}(M)}}.
    \end{equation*}
    We call it the reduced leafwise cohomology group.
    
    % \subsection{Scalar product}
    % Let $<,>_{\mathcal{F}}$ be the induced metric on $\wedge^{i}T^{*}\mathcal{F}$ from the bundle-like metric $g_{\mathcal{F}}$.
    % If we choose an orientation on $M$, it determines a volume form $\mathrm{vol}$ of $M$.
    % We set a scalar product on $\Omega^{i}_{\mathcal{F}}(M)$ by
    % \begin{equation*}
    %     (\omega,\eta)_{\mathcal{F}} = \int_{M}<\omega,\eta>_{\mathcal{F},x}\mathrm{vol}(x),
    % \end{equation*}
    % where $\omega$, $\eta$ are leafwise $i$-th forms.
    % The scalar product induces $\bar{H}^{i}_{\mathcal{F}}(M)$ a Hilbert structure.
    
    \subsection{Leafwise Hodge theorem}
    
    For a bundle-like metric $g_{\mathcal{F}}$, we have the Hodge $*$-operator.
    If we denote by $\delta$ the adjoint operator of the exterior derivative $d$, it has a decomposition into $\delta_{\mathcal{F}}$ and $\delta_{0}$
    \begin{center}
    \begin{tikzcd}
    \cdots  & \ar[l] \Omega^{i}_{\mathcal{F}}(M) \ar[d, "\oplus"] & \ar[l, "\delta_{\mathcal{F}}^{i}"] \Omega^{i+1}_{\mathcal{F}}(M)  \ar[d, "\oplus"] & \ar[l] \cdots\\
    \cdots & \ar[l]\ar[ul, ""] \Omega^{i}_{0}(M)  & \ar[l, "\delta_{\mathcal{F}}^{i}"]\ar[ul, "\delta_{0}^{i}"] \Omega^{i+1}_{0}(M) & \ar[l]\ar[ul, ""] \cdots.
    \end{tikzcd}
    \end{center}
    
    \begin{definition}[Leafwise Laplacian]
    An operator defined by
    \begin{equation*}
        \Delta_{\mathcal{F}} := d_{\mathcal{F}}\delta_{\mathcal{F}} + \delta_{\mathcal{F}}d_{\mathcal{F}} \text{ on  } \Omega_{\mathcal{F}}^{i}(M)
    \end{equation*}
    is called the \textbf{leafwise Laplacian}.
    A leafwise form $\omega\in\ker\Delta_{\mathcal{F}}$ is called a leafwise harmonic form.
    \end{definition}
    
    If we define $\Delta_{0}$ on $\Omega_{\mathcal{F}}^{i}(M)$ by $\delta_{0}d_{0}$, the restriction of the Laplacian $\Delta$ on $\Omega_{\mathcal{F}}^{i}(M)$ can be represented by
    \begin{equation*}
        \Delta|_{\Omega_{\mathcal{F}}^{i}(M)} = \Delta_{\mathcal{F}} + \Delta_{0}.
    \end{equation*}
    

    We have a significant proposition by \'Alvarez L\'opez and Kordyukov (\cite{lopez2001long}):
    \begin{proposition}[Leafwise Hodge theorem]
    Given a bundle-like metric, An leafwise cohomology class can be uniquely represented by a leafwise harmonic form.
    We have an isomorphism
    \begin{equation*}
        \bar{H}_{\mathcal{F}}^{i}(M) \cong \ker(\Delta_{\mathcal{F}}^{i}).
    \end{equation*}
    \end{proposition}